\textbackslash{}chapter\{Canonical Transformations; Poisson Brackets\}
\textbackslash{}label\{ch:5\}

coordinates. One usually finds that coordinates that simplify V make T more
complicated, and vice versa. Nevertheless, the idea of transforming to a different set of coordinates has merit. For example, a problem could be simplified
if we transformed to a new set of coordinates in which some or all of the
coordinates are ignorable. If a coordinate is ignorable, we have a partial integration (a ``first integral'') of the equation of motion because if qi is ignorable,
pi = \textbackslash{}partial L/\textbackslash{}partial $\textbackslash{}dot q$i is constant. Unfortunately, there is no known general method for
transforming to a set of variables in which the coordinates are ignorable in the
Lagrangian. Finding such a set of coordinates is a matter of intuition and luck
rather than a mathematical procedure.
On the other hand, in Hamiltonian mechanics the situation is much brighter.
Hamiltonian mechanics treats qi and pi as if both were coordinates. (In fact,
we shall see that there are transformations that convert positions into momenta
and momenta into positions.) Thus the Hamiltonian H = H(q, p, t) does not
depend on derivatives. Furthermore, the equations of motion are first order with
all of the time derivatives isolated on one side of the equation. But most importantly, Jacobi2 developed a technique for carrying out a transformation from
qi, pi to a new set of coordinates (say Qi, Pi) for which Hamilton's equations
still hold and for which the integration of the equations of motion is trivial.
Consequently, the problem of obtaining the integrated equations of motion is
reduced to the problem of finding a ``generating function'' that will yield the
desired transformation.
It should be mentioned that the procedure is complicated and you may find
some of the concepts confusing, but this should not obscure the fact that we
now have a technique for not only obtaining the equations of motion, but for
actually integrating them and determining the general solution of a dynamical
problem.3
\section*{5.2  Canonical transformations}

We have seen that the transformation equations that take us from Cartesian
coordinates to generalized coordinates are of the form
$\textbackslash{}xi$ = $\textbackslash{}xi$(q1,q2, . . . , qn; t),
i = 1, . . . , n.
3 A word on terminology. The solution of a first-order differential equation that contains as
many arbitrary constants as there are independent variables is called a ``complete integral.'' If
the solution depends on an arbitrary function, it is called the ``general integral.'' In mechanics,
one is usually interested in obtaining the complete integral.

5.2
Canonical transformations
Since such a transformation transforms the coordinates of a point from one set
of coordinates to another it is called a ``point transformation.'' (See Section 3.8.)
Point transformations can be considered to take place in configuration space.
In other words, a point transformation takes us from one set of configuration
space coordinates ($\textbackslash{}xi$) to a new set of configuration space coordinates (qi).
A point transformation yields a new set of configuration space axes.
We have also studied the Legendre transformation that allowed us to transform from the Lagrangian (a function of q, $\textbackslash{}dot q$, t ) to the Hamiltonian (a function
of q, p, t).
We now consider a particularly important type of coordinate transformations
called canonical transformations. These are phase space transformations that
take us from the set of coordinates (pi, qi) to a new set (Pi, Qi) in such a way
as to preserve the form of Hamilton's equations. Recall that the coordinates
(pi, qi) describe a system with Hamiltonian H(pi, qi, t) and are such that
$\textbackslash{}dot q$i = \textbackslash{}partial H
\textbackslash{}partial pi
,
(5.1)
and
$\textbackslash{}dot p$i = -\textbackslash{}partial H
\textbackslash{}partial qi
.
(5.2)
The ``new'' coordinates Pi and Qi can also be used to describe the Hamiltonian
of the system. The Hamiltonian in terms of the Ps and Qs will, in general,
have a different functional form than the Hamiltonian in terms of the ps and
qs. Therefore, we denote it by K(Pi, Qi, t). A canonical transformation is one
that takes us from the pi, qi to a set of coordinates Pi and Qi such that
$\textbackslash{}dot Q\_i = \textbackslash{}partial K
\textbackslash{}partial Pi
,
(5.3)
and
$\textbackslash{}dot P\_i = -\textbackslash{}partial K
\textbackslash{}partial Qi
.
(5.4)
Observe that Hamilton's equations in terms of pi, qi (Equations (5.1) and
(5.2)) and Hamilton's equations in terms of Pi, Qi (Equations (5.3) and (5.4))
have exactly the same form.
The technique for generating a canonical transformation is based on the
method used to derive Hamilton's equations. Recall that they were obtained

from a variational principle, specifically from the modified Hamilton's
principle (Equation (4.20)), which is
\textbackslash{}delta
t2
t1

pi $\textbackslash{}dot q$i -H(qi, pi, t)

dt = 0.
(5.5)
Hamilton's equations would certainly be obeyed if, in terms of the new coordinates, we had
\textbackslash{}delta
t2
t1

Pi $\textbackslash{}dot Q\_i -K(Qi, Pi, t)

dt = 0.
(5.6)
However, this is unnecessarily restrictive. For example, if F is an arbitrary
function of the coordinates and time, adding a term of the form dF
dt to the
integrand of this last expression changes nothing. This is easily appreciated as
follows: consider the expression
\textbackslash{}delta
t2
t1

Pi $\textbackslash{}dot Q\_i -K(Qi, Pi, t) + dF
dt

dt.
(5.7)
Since
\textbackslash{}delta
t2
t1
dF
dt dt = \textbackslash{}delta  (F(t2) -F(t1)) ,
and since the variation at the endpoints is required to vanish, we see that we
have just added a term equal to zero to our modified Hamilton's principle.
Therefore, nothing is changed, and Hamilton's principle still holds, but the new
form (Equation (5.7)) allows us to generate a new Hamiltonian K(Qi, Pi, t)
as well as a set of transformation equations to the new coordinates Pi and Qi.
We know that the variation expressed by Equation (5.5) is true. The variation
expressed by Equation (5.7) will certainly be true if the two integrands are
equal, that is, if
p $\textbackslash{}dot q$ -H = P $\textbackslash{}dot Q\_ -K + dF
dt .
(5.8)
(We have suppressed the summations and the subindices for simplicity.)
To see how Equation (5.8) leads to the transformation equations from p, q
to P, Q, it is convenient to assume that F depends on a particular combination
of the new and old coordinates. Thus, let us assume, for the present, that
F = F1(q, Q, t).
(5.9)

5.2
Canonical transformations
(This really means F = F1(qi, Qi, t), i = 1, . . . , n.) Note that the functional
form of F is still arbitrary. All we have done is assume that F is some function
of the new coordinates Qi and the old coordinates qi and time.
Then Equation (5.8) can be written as
p $\textbackslash{}dot q$ -H = P $\textbackslash{}dot Q\_ -K + \textbackslash{}partial F1
\textbackslash{}partial q $\textbackslash{}dot q$ + \textbackslash{}partial F1
\textbackslash{}partial Q
$\textbackslash{}dot Q\_ + \textbackslash{}partial F1
\textbackslash{}partial t .
Rearranging,

p -\textbackslash{}partial F1
\textbackslash{}partial q

$\textbackslash{}dot q$ -

P + \textbackslash{}partial F1
\textbackslash{}partial Q

$\textbackslash{}dot Q\_ = H -K + \textbackslash{}partial F1
\textbackslash{}partial t .
This equation will be satisfied if
p = \textbackslash{}partial F1
\textbackslash{}partial q ,
(5.10)
P = -\textbackslash{}partial F1
\textbackslash{}partial Q ,
(5.11)
and
K = H + \textbackslash{}partial F1
\textbackslash{}partial t .
(5.12)
At this stage it may appear we have merely stated the obvious, but actually
we have solved the problem we set out to solve: we now have an equation for
the new Hamiltonian (Equation (5.12)). Furthermore, since F1 is a function of
q and Q, then \textbackslash{}partial F1
\textbackslash{}partial q will also be a function of q and Q. So Equation (5.10) gives
us an expression of the form
p = p(q, Q, t).
This can be inverted to give
Q = Q(p, q, t).
Equation (5.11) yields
P = P(q, Q, t),
but since we know Q = Q(p, q, t) we can obtain
P = P(p, q, t).
The procedure we have used guarantees that Equations (5.1) and (5.2) are
obeyed, so the transformation is canonical.
Since this procedure is somewhat abstract, let us carry out such a transformation for a specific choice of the generating function F1(q, Q, t). For simplicity

we shall consider a system described by just two coordinates, namely p and q.
Assume that
F1(q, Q, t) = qQ.
By Equation (5.12) we see that the new Hamiltonian K is equal to the old
Hamiltonian H. By Equation (5.11) we see that the new coordinate P is
P = -\textbackslash{}partial F1
\textbackslash{}partial Q = -q.
Finally, Equation (5.10) gives
p = \textbackslash{}partial F1
\textbackslash{}partial q = Q.
Hence we see that this particular canonical transformation leads to the new set
of coordinates
P = -q,
Q = p.
In other words, the transformation changes the momentum into the coordinate and the coordinate into the momentum! This is a very good example of
the fact that the Hamiltonian formulation blurs the distinction between
momenta and coordinates, and that is why I referred to both of them simply as
``coordinates.''
The analysis we carried out assumed the generating function F was a function of the qi and Qi and we called it F1. As you might expect, useful generating functions are mixed expression involving both the old variables and
the new variables. There are four types of generating function to consider, and
these are:
F = F1(qi, Qi, t),
(5.13)
F = F2(qi, Pi, t),
F = F3(pi, Qi, t),
F = F4(pi, Pi, t).
For the three last types of generating function, the relations analogous to Equations (5.10), (5.11) and (5.12) are:

5.2
Canonical transformations
pi = \textbackslash{}partial F2
\textbackslash{}partial qi
,
(5.14)
Qi = \textbackslash{}partial F2
\textbackslash{}partial Pi
,
K = H + \textbackslash{}partial F2
\textbackslash{}partial t ,
and
qi = -\textbackslash{}partial F3
\textbackslash{}partial pi
,
(5.15)
Pi = -\textbackslash{}partial F3
\textbackslash{}partial Qi
,
K = H + \textbackslash{}partial F3
\textbackslash{}partial t ,
and
qi = -\textbackslash{}partial F4
\textbackslash{}partial pi
,
(5.16)
Qi = \textbackslash{}partial F4
\textbackslash{}partial Pi
,
K = H + \textbackslash{}partial F4
\textbackslash{}partial t .
An interesting example of the second type of generating function is
F2 =

qiPi.
(5.17)
In this case, as you can easily show, the new coordinates are related to the
old by
Qi = qi,
(5.18)
Pi = pi.
Thus, this is the identity transformation.
\begin{exercise}
Exercise 5.1 Prove that F2 =
i qiPi generates the identity transformation.
\end{exercise}

\begin{exercise}
Exercise 5.2 Obtain the transformation Equations (5.14). Hint: Subtract
PQ from F2 before plugging into Equation (5.8).
\end{exercise}

\begin{exercise}
Exercise 5.3 Let F = F4(pi, Pi, t). Determine the new Hamiltonian and
the canonical transformation equations.

\end{exercise}

\begin{example}
Example 5.1 Use a Canonical transformation to solve the harmonic oscillator problem.
Solution 5.1 The Hamiltonian for a harmonic oscillator is
H = 1
2mp2 + k
2q2.
Using the fact that the angular frequency for a harmonic oscillator is \textbackslash{}omega  =
\textbackslash{}sqrtk/m we can write the Hamiltonian in the more convenient form
H = 1
2m

p2 + m2\textbackslash{}omega 2q2
.
(5.19)
We can solve this problem by making a canonical transformation to a new
Hamiltonian K(P, Q) in which the coordinate Q is cyclic. The appropriate
generating function is
F = F1(q, Q) = m\textbackslash{}omega
2 q2 cot Q.
(I have not yet shown how to obtain the appropriate generating function, so
you will have to take this on faith.) Using Equations (5.10), (5.11) and (5.12)
we obtain
p = m\textbackslash{}omega q cot Q,
P =
m\textbackslash{}omega q2
2 sin2 Q
,
K(P, Q) = H(p, q).
A little bit of algebra yields
q2 =
m\textbackslash{}omega P sin2 Q,
and
p2 = 2m\textbackslash{}omega P cos2 Q.
These last two equations give p and q in terms of P and Q. The equation
K = H means that the new Hamiltonian and the old Hamiltonian have exactly
the same form with p and q expressed in terms of P and Q as given by the
transformation equations above. Therefore,
K(P, Q) = 2m\textbackslash{}omega P cos2 Q
2m
+
m2\textbackslash{}omega 2
m\textbackslash{}omega P sin2 Q

2m
= \textbackslash{}omega P(cos2 Q + sin2 Q) = \textbackslash{}omega P.

5.3
Poisson brackets
Thus we have obtained a Hamiltonian that is cyclic in Q, as we intended. But if
the Hamiltonian is cyclic in Q then P is constant. In this case the Hamiltonian
is the total energy E, so we can write
P = E
\textbackslash{}omega  .
Hamilton's equation for Q is
$\textbackslash{}dot Q\_ = \textbackslash{}partial K
\textbackslash{}partial P = \textbackslash{}omega .
This can be integrated immediately to yield
Q = \textbackslash{}omega t + \textbackslash{}beta ,
where \textbackslash{}beta  is a constant of integration. Transforming back to our original coordinates
q =

2P
m\textbackslash{}omega  sin Q,
so
q =

2E
m\textbackslash{}omega 2 sin(\textbackslash{}omega t + \textbackslash{}beta ),
and the problem is solved.
This application for solving a simple problem has been likened to using a
sledgehammer to crack a peanut.4 However, it emphasizes that in mechanics
the Hamiltonian is a theoretical rather than a practical tool. (Of course, this is
not true when you are doing Quantum Mechanics where the use of the Hamiltonian is the only reasonable way to solve many problems.)
\end{example}

\begin{exercise}
Exercise 5.4 Show
that,
for
the
harmonic
oscillator,
p =
\textbackslash{}sqrt
2mE
cos(\textbackslash{}omega t + \textbackslash{}beta ).
\end{exercise}

\section*{5.3  Poisson brackets}

We now introduce the Poisson brackets. The notation used may be a bit
confusing at first, but you will probably quickly appreciate how useful the
Poisson brackets are in the development of Hamiltonian dynamics. They also
4 Herbert Goldstein, Classical Mechanics, Addison-Wesley Pub. Co., Reading MA, USA, 1950,
page 247.

pave the way for transforming classical dynamics into quantum mechanics.
On a practical level, the Poisson brackets give us an easy way to determine
whether or not a transformation from one set of variables to another is
canonical.
Let p and q be canonical variables, and let u and v be functions of p and q.
The Poisson bracket of u and v is defined as
[u, v]p,q \textbackslash{}equiv \textbackslash{}partial u
\textbackslash{}partial q
\textbackslash{}partial v
\textbackslash{}partial p -\textbackslash{}partial u
\textbackslash{}partial p
\textbackslash{}partial v
\textbackslash{}partial q .
(5.20)
Generalizing to a system of n degrees of freedom we have
[u, v] =
n

i=1
\textbackslash{}partial u
\textbackslash{}partial qi
\textbackslash{}partial v
\textbackslash{}partial pi
-\textbackslash{}partial u
\textbackslash{}partial pi
\textbackslash{}partial v
\textbackslash{}partial qi

.
Using the Einstein summation convention this is written
[u, v] = \textbackslash{}partial u
\textbackslash{}partial qi
\textbackslash{}partial v
\textbackslash{}partial pi
-\textbackslash{}partial u
\textbackslash{}partial pi
\textbackslash{}partial v
\textbackslash{}partial qi
.
(5.21)
From the definition of the Poisson bracket it is obvious that
[qi, qj] = [pi, pj] = 0,
(5.22)
and that
[qi, pj] = -[pi, qj] = \textbackslash{}delta ij.
(5.23)
It is interesting to note that relationships involving Poisson brackets can
be transformed into quantum mechanical relationships by replacing the
Poisson bracket with the quantum mechanical commutator using the simple
prescription
[u, v] \textbackslash{}to 1
iℏ(ˆuˆv -ˆv ˆu),
(5.24)
where u, v are classical functions and ˆu, ˆv are the corresponding quantum
mechanical operators.
A Poisson bracket is invariant under a change in canonical variables.
That is,
[u, v]p,q = [u, v]P,Q.
In other words, Poisson brackets are canonical invariants. This relationship
gives us an easy way to determine whether or not a set of variables is canonical.

5.4
The equations of motion in terms of Poisson brackets
Rules for manipulating Poisson brackets
There are a few simple rules involving Poisson brackets that can (mostly) be
proved immediately by writing out the definition of the Poisson brackets. These
rules essentially define the arithmetic of Poisson brackets. They are:
[u, u] = 0
(5.25)
[u, v] = -[v, u]
(5.26)
[au + bv, w] = a[u, w] + b[v, w]
(5.27)
[uv, w] = [u, w]v + u[v, w].
(5.28)
A useful relationship which is somewhat harder to prove is called Jacobi's
identity. It is:
[u, [v, w]] + [v, [w, u]] + [w, [u, v]] = 0.
(5.29)
\begin{exercise}
Exercise 5.5 Prove Equations (5.22) and (5.23).
\end{exercise}

\begin{exercise}
Exercise 5.6 Prove that Poisson brackets are invariant under canonical
transformations.
\end{exercise}

\begin{exercise}
Exercise 5.7 Prove Equations (5.25) and (5.26).
\end{exercise}

\section*{5.4  The equations of motion in terms of Poisson brackets}

We have expressed the equations of motion for a mechanical system in a variety
of ways, including Newton's second law, the Lagrange equations, and Hamilton's equations. The equations of motion can also be expressed in terms of the
Poisson brackets, as we now demonstrate.
Recall that Hamilton's equations are a set of 2n first-order equations for
$\textbackslash{}dot q$i and $\textbackslash{}dot p$i, that is, for the first derivative of the variables. (This differs from
the Lagrange equations or Newton's laws which are a set of n second-order
equations for $\textbackslash{}ddot q$i.)
If u is a function of qi, pi, t, we can write
du
dt = \textbackslash{}partial u
\textbackslash{}partial qi
$\textbackslash{}dot q$i + \textbackslash{}partial u
\textbackslash{}partial pi
$\textbackslash{}dot p$i + \textbackslash{}partial u
\textbackslash{}partial t .

But $\textbackslash{}dot q$i and $\textbackslash{}dot p$i can be expressed using Hamilton's equations so we obtain
du
dt = \textbackslash{}partial u
\textbackslash{}partial qi
\textbackslash{}partial H
\textbackslash{}partial pi
-\textbackslash{}partial u
\textbackslash{}partial pi
\textbackslash{}partial H
\textbackslash{}partial qi
+ \textbackslash{}partial u
\textbackslash{}partial t ,
= [u, H] + \textbackslash{}partial u
\textbackslash{}partial t .
(5.30)
If u = constant, du
dt = 0 and [u, H] = -\textbackslash{}partial u
\textbackslash{}partial t . Therefore, if u does not depend
explicitly on t, then [u, H] = 0.5
Now suppose that u = q. We have
$\textbackslash{}dot q$ = [q, H].
(5.31)
Similarly,
$\textbackslash{}dot p$ = [p, H].
(5.32)
Thus we have written Hamilton's equations of motion in terms of Poisson
brackets. Furthermore, if we let u = H we have
dH
dt = [H, H] + \textbackslash{}partial H
\textbackslash{}partial t ,
or
dH
dt = \textbackslash{}partial H
\textbackslash{}partial t ,
(5.33)
as we have already seen (Equation (4.19)).
In terms of symplectic notation where the column matrix \textbackslash{}eta  has components
q1, . . . , pn, we can combine Equations (5.31) and (5.32) into one, thus
$\textbackslash{}dot\textbackslash{}$eta  = [\textbackslash{}eta , H],
(5.34)
which is a prescription for generating the values of q1, . . . , pn at some later
time if their values at time t are known, because
\textbackslash{}eta (t + dt) = \textbackslash{}eta (t) + $\textbackslash{}dot\textbackslash{}$eta dt.
5.4.1 Infinitesimal canonical transformations
In an infinitesimal canonical transformation the new coordinates differ
only infinitesimally from the old coordinates. Therefore, the transformation
equations from pi, qi to Pi, Qi will have the form
5 In quantum mechanical terminology we would say that ˆu commutes with ˆH. See
Equation (5.24). You probably recall from your quantum mechanics course that conserved
quantities commute with the Hamiltonian.

5.4
The equations of motion in terms of Poisson brackets
Qi = qi + \textbackslash{}delta qi,
Pi = pi + \textbackslash{}delta pi.
(A word on the notation. Here \textbackslash{}delta qi and \textbackslash{}delta pi are infinitesimal changes in the qi
and pi and are not virtual displacements.)
For the identity transformation we wrote
pi = \textbackslash{}partial F2
\textbackslash{}partial qi
,
Qi = \textbackslash{}partial F2
\textbackslash{}partial Pi
,
with F2 =
qiPi. (See Equations (5.14).) Therefore, the infinitesimal canonical transformation which differs infinitesimally from the identity transformation has a ``generating function'' given by
F2(q, P, t) =
qiPi + ϵG(q, P, t),
where ϵ is an infinitesimal quantity and G is an arbitrary function. Then, by
the rules for F2 transformations (Equations (5.14)) we see that
pj = \textbackslash{}partial F2
\textbackslash{}partial qj
= Pj + ϵ \textbackslash{}partial G
\textbackslash{}partial qj
,
Qj = \textbackslash{}partial F2
\textbackslash{}partial Pj
= qj + ϵ \textbackslash{}partial G
\textbackslash{}partial Pj
.
Now,
\textbackslash{}delta qj = Qj -qj = ϵ \textbackslash{}partial G
\textbackslash{}partial Pj
,
and to first order in ϵ we can write
\textbackslash{}delta qj = ϵ \textbackslash{}partial G
\textbackslash{}partial pj
.
(5.35)
Similarly, \textbackslash{}delta pj = Pj -pj so
\textbackslash{}delta pj = -ϵ \textbackslash{}partial G
\textbackslash{}partial qj
.
(5.36)
In symplectic notation Equations (5.35) and (5.36) are expressed as
\textbackslash{}delta \textbackslash{}eta  = ϵ [\textbackslash{}eta , G] .
(5.37)

Since F2 is guaranteed to generate a canonical transformation, we appreciate
that (5.37) is indeed an infinitesimal canonical transformation.
In general
\textbackslash{}zeta  = \textbackslash{}eta  + \textbackslash{}delta \textbackslash{}eta ,
(5.38)
is a transformation from q1, . . . , pn to a nearby point in phase space
q1 + \textbackslash{}delta q1, . . . , pn + \textbackslash{}delta pn. But if the quantity G is taken to be the Hamiltonian H, the transformation is a displacement in time from q1(t), . . . , pn(t) to
q1(t + dt), . . . , pn(t + dt). (This canonical transformation in time is sometimes called a ``contact'' transformation.) Thus, the Hamiltonian is the generator of the motion of the system in time; it generates the canonical variables at
a later time.
We now consider the change in a function u = u(q, p) under a canonical
transformation. Note that there are two ways to interpret a canonical transformation. On the one hand, it changes the description of a system from an
old set of canonical variables q, p to a new set of canonical variables Q, P.
Under such a transformation a function u may have a different form but it
will still have the same value. Thus, if q1(t), . . . , pn(t) are denoted by A and
Q1(t), . . . , Pn(t) are denoted by A′, the transformed function u(A′) has the
same value as the original function u(A):
u(A′) = u(A).
(The value of the total angular momentum of a rotating body is the same
regardless of the set of coordinates used to evaluate it.) Of course, u(A′) will
in general have a different mathematical form than u(A). (In Cartesian coordinates the angular momentum of a mass point moving in a circle is m(x $\textbackslash{}dot y$ -y $\textbackslash{}dot x$)
and in polar coordinates it is mr2 $\textbackslash{}dot\textbackslash{}$theta .)
On the other hand, if we consider an infinitesimal contact transformation
that generates a translation in time, the interpretation is quite different. Now
the transformation takes us from A, representing the values q1(t), . . . , pn(t),
to B, representing q1(t + dt), . . . , pn(t + dt), in the same phase space (with
the same set of phase space axes) but at a later time. Since we are in the same
phase space and using the same set of coordinates, the form of u is conserved
but its value changes. Let us represent such a change in the function u by
\textbackslash{}partial u = u(B) -u(A).

5.4
The equations of motion in terms of Poisson brackets
But since u = u(q, p) we write
\textbackslash{}partial u = \textbackslash{}partial u
\textbackslash{}partial qi
\textbackslash{}delta qi + \textbackslash{}partial u
\textbackslash{}partial pi
\textbackslash{}delta pi
= \textbackslash{}partial u
\textbackslash{}partial qi
ϵ \textbackslash{}partial G
\textbackslash{}partial pi + \textbackslash{}partial u
\textbackslash{}partial pi

-ϵ \textbackslash{}partial G
\textbackslash{}partial qi

= ϵ

\textbackslash{}partial u
\textbackslash{}partial qi
\textbackslash{}partial G
\textbackslash{}partial pi
-\textbackslash{}partial u
\textbackslash{}partial pi
\textbackslash{}partial G
\textbackslash{}partial qi

,
or
\textbackslash{}partial u = ϵ[u, G].
(5.39)
Now let us consider a change in the Hamiltonian under a contact transformation. The Hamiltonian differs from an ordinary function u in the following way.
Under a canonical transformation u(A) \textbackslash{}to u(A′) the value of u is constant, but
for the Hamiltonian, a canonical transformation from A to A′ can result in a
function with a different value. In fact, the transformed Hamiltonian will generally be a completely different function. This is because the Hamiltonian is
not a function having a specific value, but rather a function which defines the
canonical equations of motion. Therefore, we must write
H(A) \textbackslash{}to K(A′).
As we know,
K = H + \textbackslash{}partial F
\textbackslash{}partial t .
But for an infinitesimal canonical transformation, we want K to differ from H
by only an infinitesimal amount, so F is approximately equal to the identity
transformation. As before, we write
F2 =

qiPi + ϵG(q, P, t).
Then,
K(A′) = H(A) + \textbackslash{}partial
\textbackslash{}partial t

qiPi + ϵG(q, P, t)

= H(A) + ϵ \textbackslash{}partial G
\textbackslash{}partial t ,
and
\textbackslash{}partial H = H(B) -K(A′) = H(B) -H(A) -ϵ \textbackslash{}partial G
\textbackslash{}partial t .
But
H(B) -H(A) = ϵ[H, G],

so
\textbackslash{}partial H = ϵ[H, G] -ϵ \textbackslash{}partial G
\textbackslash{}partial t .
Now by Equation (5.30),
dG
dt = [G, H] + \textbackslash{}partial G
\textbackslash{}partial t .
Therefore
ϵ \textbackslash{}partial G
\textbackslash{}partial t = ϵ dG
dt -ϵ[G, H],
and
\textbackslash{}partial H = ϵ[H, G] -ϵ \textbackslash{}partial G
\textbackslash{}partial t + ϵ[G, H] = -ϵ \textbackslash{}partial G
\textbackslash{}partial t .
(5.40)
If G is a constant of the motion, dG
dt = 0 and \textbackslash{}partial H = 0. That is, if the generating function is a constant of the motion, the Hamiltonian is invariant. This
is closely related to the relationship between symmetries and constants of the
motion. Recall that under a translation involving an ignorable coordinate, the
Hamiltonian is invariant and the conjugate generalized momentum is constant.
But if the Hamiltonian is invariant, then G, the generator of the infinitesimal canonical transformation, must be a constant of the motion (because then
dG
dt = 0). That is, the symmetry (the cyclic nature of one of the coordinates)
implies a constant of the motion (G).
5.4.2 Canonical invariants
So far we have considered two quantities that are canonical invariants, that is,
quantities that do not change when the variables undergo a canonical transformation from the p, q to the P, Q. The first set of such canonical invariants is,
of course, Hamilton's equations. The second is the Poisson bracket. Any equation expressed in terms of Poisson brackets will be invariant under a canonical
transformation of variables. One can actually develop a system of mechanics
based on Poisson brackets in which all equations are the same for any set of
canonical variables.
A third quantity that can be shown to be a canonical invariant is the volume
element in phase space. In other words, if the volume element in terms of the
variables qi, pi is
d\textbackslash{}eta  = dq1dq2 \textbackslash{}cdot  \textbackslash{}cdot  \textbackslash{}cdot  dqn dp1dp2 \textbackslash{}cdot  \textbackslash{}cdot  \textbackslash{}cdot  dpn,

5.4
The equations of motion in terms of Poisson brackets
and the volume element in terms of the canonical variables Qi, Pi is
d\textbackslash{}zeta  = dQ1dQ2 \textbackslash{}cdot  \textbackslash{}cdot  \textbackslash{}cdot  dQn dP1dP2 \textbackslash{}cdot  \textbackslash{}cdot  \textbackslash{}cdot  dPn,
then
d\textbackslash{}eta  = d\textbackslash{}zeta .
(5.41)
Let us consider the proof of this assertion. The proof is based on the fact that
the transformation of the infinitesimal volume element expressed in one set of
variables (x1, x2, . . . , xn) to another set (q1, q2, . . . , qm) is given by
dx1dx2 \textbackslash{}cdot  \textbackslash{}cdot  \textbackslash{}cdot  dxn = Ddq1dq2 \textbackslash{}cdot  \textbackslash{}cdot  \textbackslash{}cdot  dqm,
where D is the Jacobian determinant defined by
D = \textbackslash{}partial (q1, . . . , qm)
\textbackslash{}partial (x1, . . . , xn) =

\textbackslash{}partial q1
\textbackslash{}partial x1
\textbackslash{}partial q1
\textbackslash{}partial x2
. . .
\textbackslash{}partial q1
\textbackslash{}partial xn
...
\textbackslash{}partial qm
\textbackslash{}partial x1
\textbackslash{}partial qm
\textbackslash{}partial x2
\textbackslash{}cdot  \textbackslash{}cdot  \textbackslash{}cdot
\textbackslash{}partial qm
\textbackslash{}partial xn

.
The volume of some region of phase space expressed in terms of the two sets
of canonical variables (q1, . . . , pn) and (Q1, . . . , Pn) is

. . .

dQ1dQ2 \textbackslash{}cdot  \textbackslash{}cdot  \textbackslash{}cdot  dPn =

. . .

Ddq1dq2 \textbackslash{}cdot  \textbackslash{}cdot  \textbackslash{}cdot  dpn.
The volume is invariant if and only if D = 1. Thus the proof of the invariance
of the volume under a canonical transformation is reduced to proving that the
Jacobian determinant is equal to unity.
Note that
D = \textbackslash{}partial (Q1, . . . , Pn)
\textbackslash{}partial (q1, . . . , pn) =

\textbackslash{}partial Q1
\textbackslash{}partial q1
\textbackslash{}partial Q1
\textbackslash{}partial q2
\textbackslash{}cdot  \textbackslash{}cdot  \textbackslash{}cdot
\textbackslash{}partial Q1
\textbackslash{}partial pn
...
\textbackslash{}partial Pn
\textbackslash{}partial q1
\textbackslash{}partial Pn
\textbackslash{}partial q2
\textbackslash{}cdot  \textbackslash{}cdot  \textbackslash{}cdot
\textbackslash{}partial Pn
\textbackslash{}partial pn

.
Now the Jacobian can be treated somewhat like a fraction and the value of D
is not changed if we divide top and bottom by \textbackslash{}partial (q1, . . . , qn, P1, . . . , Pn), thus
D = \textbackslash{}partial (Q1, . . . , Pn)
\textbackslash{}partial (q1, . . . , pn) =
\textbackslash{}partial (Q1,...,Pn)
\textbackslash{}partial (q1,..., qn,P1,...,Pn)
\textbackslash{}partial (q1,...,pn)
\textbackslash{}partial (q1,..., qn,P1,...,Pn)
.
Another property of Jacobians (which follows from the properties of
determinants) is that if the same quantities appear in the numerator and

denominator, they can be dropped, yielding a Jacobian of lower order.6
Therefore we can write
D = \textbackslash{}partial (Q1, . . . , Qn)/\textbackslash{}partial (q1, . . . , qn)
\textbackslash{}partial (p1, . . . , pn)/\textbackslash{}partial (P1, . . . , Pn) = \textbar{}n\textbar{}
\textbar{}d\textbar{},
where we dropped the Ps in the numerator and the qs in the denominator.
(The quantity \textbar{}n\textbar{} represents the determinant of the numerator and \textbar{}d\textbar{} represents the determinant of the denominator.) The ikth element of the determinant
in the numerator is
nik = \textbackslash{}partial Qi
\textbackslash{}partial qk
,
and the kith element of the denominator is
dki = \textbackslash{}partial pk
\textbackslash{}partial Pi
.
If the generating function of the canonical transformation has the form
F = F2(qi, Pi, t), then by Equations (5.14)
pi = \textbackslash{}partial F
\textbackslash{}partial qi
,
Qi = \textbackslash{}partial F
\textbackslash{}partial Pi
,
and consequently the ikth element of the numerator is
nik =
\textbackslash{}partial
\textbackslash{}partial qk
\textbackslash{}partial F
\textbackslash{}partial Pi
=
\textbackslash{}partial 2F
\textbackslash{}partial qk\textbackslash{}partial Pi
,
and the kith element of the denominator is
dki =
\textbackslash{}partial
\textbackslash{}partial Pi
\textbackslash{}partial F
\textbackslash{}partial qk
=
\textbackslash{}partial 2F
\textbackslash{}partial qk\textbackslash{}partial Pi
,
so the two determinants differ only in the fact that the rows and columns are
interchanged. This does not affect the value of a determinant, so we conclude
that D = 1 and the assertion is proved.
\begin{exercise}
Exercise 5.8 Show that for the transformation from Cartesian to polar
coordinates that D = r.
6 See, for example, Section 6.4.4 of Mathematical Methods for Physics and Engineering by
K. F. Riley, M. P. Hobson and S. J. Bence, 2nd Edn, Cambridge University Press, 2002.

5.4
The equations of motion in terms of Poisson brackets
\end{exercise}

\begin{exercise}
Exercise 5.9 Consider a canonical transformation from (q, p) to (Q, P).
(Here n = 1.) (a) Show that
\textbackslash{}partial (Q, P)
\textbackslash{}partial (q, p) = \textbackslash{}partial (Q, P)/\textbackslash{}partial (q, P)
\textbackslash{}partial (q, p)/\textbackslash{}partial (q, P) .
(b) Show that
\textbackslash{}partial (Q, P)
\textbackslash{}partial (q, P) = \textbackslash{}partial Q
\textbackslash{}partial q .
5.4.3 Liouville's theorem
Recall the analogy between the phase fluid and real fluids and consider that
some fluids are incompressible. (Water is a reasonably good example of an
incompressible fluid.) From fluid dynamics it is known that if a fluid is incompressible, then the velocity field v = v(x, y, z) has zero divergence: \textbackslash{}nabla\textbackslash{}cdot  v = 0.
The phase fluid behaves like a 2n-dimensional incompressible fluid. The
``velocity'' components for this fluid are $\textbackslash{}dot q$i and $\textbackslash{}dot p$i. Therefore, the condition
for incompressibility is
n

i=1
\textbackslash{}partial $\textbackslash{}dot q$i
\textbackslash{}partial qi
+ \textbackslash{}partial $\textbackslash{}dot p$i
\textbackslash{}partial pi

= 0.
(5.42)
Applying the divergence theorem to the velocity of an incompressible fluid we
find that the flux of the fluid across any closed surface is zero. Similarly, the
flux of the phase fluid over any closed surface in phase space is zero. This fact,
which was discovered by Liouville, is called ``Liouville's theorem.'' Since the
development of a system in time can be treated as a canonical transformation,
this is just another way of stating that the volume of a region of the phase fluid
is constant in time, and we have just proved that this is a property of canonical
transformations.
Liouville's theorem is particularly useful in statistical mechanics. For example, in considering a gas the system will probably contain about 1026 molecules.
It is clearly out of the question to follow the motion of each individual particle. In statistical mechanics this type of problem is dealt with by imagining
an ensemble of such systems and evaluating ensemble averages of the parameters of interest. For example, the ensemble could consist of a large number of
identical systems of gas molecules differing only in their initial conditions. In
phase space, each member of the ensemble is represented by a single point and

the whole ensemble is represented by a large number of points. Liouville's theorem states that the density of these points remains constant in time. Note that
if the phase fluid is incompressible, it has a constant density. By the constant
density of the phase fluid we mean that the number of points in a given volume
does not change. Each point represents a system. As time goes on, the points
in a volume move along their phase paths and the surface enclosing them will
be distorted, but it will have the same volume as before. Now the volume of a
region of phase space is given by
\textbackslash{}tau  =

dq1 \textbackslash{}cdot  \textbackslash{}cdot  \textbackslash{}cdot  dqndp1 \textbackslash{}cdot  \textbackslash{}cdot  \textbackslash{}cdot  dpn.
Liouville's theorem gives us an additional constant of the motion, namely
\textbackslash{}tau  = constant.
Another point of view is to consider an infinitesimal region of phase space
containing a number of systems. The density is just the number of systems
divided by the volume. At a later time, these systems will have moved to a different region of phase space. The region enclosing the points will have changed
shape, but will have the same volume as before because the volume element
in phase space is a canonical invariant. But if the number of systems is constant and the volume is constant, then the density is constant and the theorem
is proved.
It is interesting to note that if the density is denoted by \textbackslash{}rho  then in general,
d\textbackslash{}rho
dt = [\textbackslash{}rho , H] + \textbackslash{}partial \textbackslash{}rho
\textbackslash{}partial t .
But since d\textbackslash{}rho
dt = 0 we have
\textbackslash{}partial \textbackslash{}rho
\textbackslash{}partial t = -[\textbackslash{}rho , H].
\end{exercise}

\begin{exercise}
Exercise 5.10 Show that if \textbackslash{}nabla\textbackslash{}cdot  v = 0 then there is no net flow of material across the boundary of a given volume. (Hint: Use the divergence
theorem.)
\end{exercise}

\begin{exercise}
Exercise 5.11 Show that Equation (5.42) is true in general.
5.4.4 Angular momentum
Another application of infinitesimal canonical transformations involves the
angular momentum of a system. If a system is symmetric with respect to a

5.5
Angular momentum in Poisson brackets
rotation, then the rotation angle (say \textbackslash{}theta ) is cyclic and the conjugate generalized
momentum (the angular momentum L) will be constant.7 Conversely, if the
generating function is taken to be the angular momentum, the transformation
is a rigid rotation of the system. This can be shown as follows. Let
G = Lz =

(ri \textbackslash{}times  pi)z =

(xipyi -yipxi).
(We selected the z component of angular momentum for simplicity; this does
not imply any lack of generality.) If we carry out an infinitesimal rotation d\textbackslash{}theta
about the z axis, it is easy to show that
\textbackslash{}delta $\textbackslash{}xi$ = -yid\textbackslash{}theta
\textbackslash{}delta yi = xid\textbackslash{}theta
\textbackslash{}delta zi = 0,
and
\textbackslash{}delta pxi = -pyid\textbackslash{}theta
\textbackslash{}delta pyi = pxid\textbackslash{}theta
\textbackslash{}delta pzi = 0.
Now we have seen that \textbackslash{}delta qj = ϵ \textbackslash{}partial G
\textbackslash{}partial pj (Equation (5.35)) and that \textbackslash{}delta pj = -ϵ \textbackslash{}partial G
\textbackslash{}partial qj
(Equation (5.36)), so if we use d\textbackslash{}theta  for the infinitesimal parameter ϵ,
\textbackslash{}delta $\textbackslash{}xi$ = ϵ \textbackslash{}partial G
\textbackslash{}partial pxi
= ϵ
\textbackslash{}partial
\textbackslash{}partial pxi

xjpyj -yjpxj

= ϵ(-yi) = -yid\textbackslash{}theta ,
as expected. The other relations follow similarly.
\end{exercise}

\section*{5.5  Angular momentum in Poisson brackets}

We have seen that the infinitesimal canonical transformation generates a change
in a function u given by
\textbackslash{}partial u = ϵ[u, G].
In the ``active'' interpretation we think of the ICT as a transformation that generates the (new) value of the function due to some change in the system parameters. In particular, if G is the Hamiltonian, then we let ϵ = dt and \textbackslash{}partial u is the
change in the function u as the parameter time goes from t to t + dt. If G is
some other function, then \textbackslash{}partial u represents a different sort of change in u.
7 Do not confuse the vector angular momentum L with the Lagrangian.

If u is a particular component of a vector F, say Fi, then the change in Fi
due to an infinitesimal canonical transformation is
\textbackslash{}partial Fi = ϵ[Fi, G].
If G is a component of the angular momentum (G = ˆn \textbackslash{}cdot  L) and ϵ is taken to be
d\textbackslash{}theta , then the ICT gives the change in Fi due to a rotation of the system through
the angle d\textbackslash{}theta  about ˆn:
\textbackslash{}partial Fi = d\textbackslash{}theta [Fi, ˆn \textbackslash{}cdot  L].
Note that this ICT gives the change in the component Fi along some direction fixed in space, that is, along a fixed axis. For example, Fi might be F \textbackslash{}cdot  \textbackslash{}hat\{k\}
where the Cartesian unit vectors î, \textbackslash{}hat\{j\}, \textbackslash{}hat\{k\} are assumed to be fixed in space and not
affected by the rotation. Clearly, the components of some vectors are affected
by the rotation and others are not. Axes fixed in the body rotate with it; axes
fixed in space do not. An external vector, such as the gravitational force, is
obviously not affected by a rotation of the body. On the other hand, a component of the angular momentum of the system would, in general, be affected by
such a rotation. A ``system vector'' is a vector whose components depend on
the orientation of the body and which is affected by a rotation of the system.
The change in a system vector due to an infinitesimal rotation d\textbackslash{}theta  about some
axis ˆn is, according to vector analysis,
dF = ˆnd\textbackslash{}theta \textbackslash{}times F,
whereas the Poisson bracket formulation yields
\textbackslash{}partial F = d\textbackslash{}theta [F, ˆn \textbackslash{}cdot  L].
The two formulations must be equivalent, so
[F, ˆn\textbackslash{}cdot L] = ˆn \textbackslash{}times  F.
(5.43)
This equation is particularly interesting and useful because all reference to the
rotation has been deleted and is valid for all system vectors. It gives the relationship between the Poisson bracket of the system vector with a component of
the angular momentum [F, ˆn\textbackslash{}cdot L] and the cross product of the unit vector in the
direction of the angular momentum component and the system vector ˆn\textbackslash{}times F.
Sometimes this last relationship is more easily used in terms of components in
which case it reads

5.5
Angular momentum in Poisson brackets
[Fi, Lj] = ϵijkFk,
(5.44)
where ϵijk is the Levi--Civita density.8
Another useful relationship involves the Poisson bracket of the dot product
of two system vectors (say F and V) and a component of L. Using (5.43):
[F \textbackslash{}cdot  V, ˆn \textbackslash{}cdot  L] = F \textbackslash{}cdot  [V, ˆn \textbackslash{}cdot  L] + V \textbackslash{}cdot  [F, ˆn \textbackslash{}cdot  L]
(5.45)
= F \textbackslash{}cdot  (ˆn \textbackslash{}times  V) + V \textbackslash{}cdot  (ˆn \textbackslash{}times  F)
= (F \textbackslash{}times  ˆn) \textbackslash{}cdot  V + V \textbackslash{}cdot  (ˆn \textbackslash{}times  F) = 0,
where we exchanged the dot and cross on the first term. Now if F and V are
both taken to be L, Equation (5.44) gives the following relation between the
components of the angular momentum
[Li, Lj] = ϵijkLk,
(5.46)
and Equation (5.45) gives
[L2, ˆn \textbackslash{}cdot  L] = 0.
(5.47)
It can easily be shown that the Poisson bracket of any two constants of the
motion is also a constant of the motion. (This is called Poisson's theorem.)
Therefore, if Lx and Ly are constants of the motion, then by (5.46), so too is
Lz. Recall that according to Equation (5.22) the Poisson bracket of any two
canonical momentum components is zero. But [Li, Lj] = ϵijkLk, so Li and
Lj cannot be canonical variables. More generally, no two components of L
can simultaneously be canonical variables. One the other hand, the relation
[L2, ˆn \textbackslash{}cdot  L] = 0 shows that the magnitude of L (or L2) and a component of L
can simultaneously be canonical variables.9
Furthermore if p is the linear momentum of the system and if pz is constant,
then so too are px and py. This is easily proved because by Equation (5.44)
[pz, Lx] = py,
and, according to Poisson's theorem, if pz and Lx are constants, then py is also
a constant. Similarly,
[pz, Ly] = -px,
8 The Levi-Civita density ϵijk is defined to be zero if any two indices are equal; it is +1 if ijk
are even permutations of 123, and it is -1 if ijk are odd permutations of 123.
9 This statement carries over into quantum mechanics where we find that no two components of
L can simultaneously have eigenvalues, but, for example, Lz and L2 can be quantized
together.

and so px is also constant. Thus, both L and p are conserved if Lx, Ly and pz
are constant.
All of these relationships are probably familiar to you from your study of
quantum mechanics.
\section*{5.6  Problems}

5.1 (a) Prove that the transformation
Q = p/ tan q,
P = log(sin q/p),
is canonical. (b) Find the generating function F1(q, Q) for this transformation.
5.2 Assume H(q1, q2, p1, p2) = q1p1 -q2p2 -aq2
1 + bq2
2, where a and b
are constants. Show that the time derivative of q1q2 is zero. (Use Poisson
brackets.)
5.3 Under what condition is the following transformation not canonical?
Q = q + ip,
P = Q∗.
\section*{5.4  Consider the generating function}

F2(q, P) = (q + P)2.
Obtain the transformation equations.
\section*{5.5  The Hamiltonian for a falling body is}

H = p2
2m + mgz,
where g is the acceleration of gravity. Determine the generating function
F4(p, P) if K = P.
5.6 Does F1(q, Q) = q2 + Q4 generate a canonical transformation?
5.7 Suppose f and g are constants of the motion. Prove that their Poisson
bracket is also an integral of the motion.
5.8 Prove Jacobi's identity (Equation (5.29)). Show that the commutator
[A, B] = AB -BA also satisfies a relation similar to the Jacobi identity. That is, show that [A, BC] = [A, B]C + B[A, C].
5.9 Consider a generating function of the form F3(pi, Qi, t). (a) Derive
expressions for qi, Pi and K, as given by Equations (5.15). (b) Determine
q, P, and K, if F3 = pQ.

5.6
Problems
5.10 Assume q and p are canonical. Show that if Q and P are canonical, then
the following conditions hold
\textbackslash{}partial Q
\textbackslash{}partial q = \textbackslash{}partial p
\textbackslash{}partial P , \textbackslash{}partial Q
\textbackslash{}partial p = -\textbackslash{}partial q
\textbackslash{}partial P , \textbackslash{}partial P
\textbackslash{}partial q = -\textbackslash{}partial p
\textbackslash{}partial Q, \textbackslash{}partial P
\textbackslash{}partial p = \textbackslash{}partial q
\textbackslash{}partial Q.
Prove directly and by using Poisson brackets.

Hamilton--Jacobi theory
In the previous chapter it was mentioned that there is no general technique for
solving the n coupled second-order Lagrange equations of motion, but that
Jacobi had derived a general method for solving the 2n coupled canonical
equations of motion, allowing one to determine all the position and momentum
variables in terms of their initial values and the time.
There are two slightly different ways to solve Hamilton's canonical equations. One is more general, whereas the other is a bit simpler, but is only valid
for systems in which energy is conserved. We will go through the procedure
for the more general method, then solve the harmonic oscillator problem by
using the second method.
Both methods involve solving a partial differential equation for the quantity
S that is called ``Hamilton's principal function.'' The problem of solving the
entire system of equations of motion is reduced to solving a single partial differential equation for the function S. This partial differential equation is called
the ``Hamilton--Jacobi equation.'' Reducing the dynamical problem to solving
just one equation is quite satisfying from a theoretical point of view, but it is
not of much help from a practical point of view because the partial differential equation for S is often very difficult to solve. Problems that can be solved
by obtaining the solution for S can usually be solved more easily by other
means.
Nevertheless, the solution of the Hamilton--Jacobi equation is important for
a variety of reasons. For one thing, the analysis allows one to draw an ``opticalmechanical analogy'' between optical rays and mechanical phase space paths.
That is, we can treat waves from a mechanical point of view. This led, of
course, to ``wave mechanics'' and quantum theory.